% =========================================================================
% model_section_v3.tex — Two-channel comparative statics + (S,s) extension
%
% DRAFT FOR PI LINE-BY-LINE VERIFICATION. Written 2026-07-07 to be \input
% after "The Rust Model Applied to Codebase Decisions" in DRAFT_paper.tex
% (notation matched: q_t state, a_t in {0,1}, c(q,theta_c)=theta_1 q +
% theta_2 q^2, kappa rationalization cost, reset to q_0 + eta, discount
% beta, T1EV shocks). New parameter: alpha = AI-assistance intensity.
%
% Integration notes:
%  - Propositions use \newtheorem{proposition}{Proposition} and
%    \newtheorem{remark}{Remark} — add to the preamble if not present.
%  - Appendix proofs are at the bottom of this file inside
%    \section*{Appendix ...}; move to the paper's appendix.
%  - NEW CITATIONS in model_v3_new_refs.bib are UNVERIFIED — per project
%    citation policy, verify each against the primary source before merging.
%  - FORWARD REFERENCES that must exist in the rewritten draft:
%    \label{sec:results}, \label{subsec:adoption_design},
%    \label{fig:event_study_components} (11_event_study_components.png).
%    \label{subsec:events} already exists in DRAFT_paper.tex.
% =========================================================================

\subsection{Comparative Statics: Two Channels of AI Adoption}
\label{subsec:twochannel}

AI-assisted development enters the model through two distinct structural
objects. Let $\alpha \geq 0$ index the intensity of AI assistance in the
firm's development process (Section~\ref{subsec:events} describes its
empirical counterpart: configuration-artifact adoption dates and the share
of AI co-authored commits).

\begin{enumerate}
    \item \textbf{Cost channel:} the rationalization cost falls,
    $\kappa = \kappa(\alpha)$ with $\kappa'(\alpha) < 0$. AI agents execute
    large, mechanical rewrites --- migrations, API-surface changes,
    dependency upgrades --- at a fraction of their pre-AI engineering cost.
    \item \textbf{Transition channel:} the codebase degrades faster under
    continuation. Write the continuation transition as
    $q_{t+1} \sim G_\alpha(\cdot \mid q_t)$, the distribution of
    $q_t + \delta_d(q_t, \eta_t; \alpha)$, with $G_\alpha$ shifting toward
    faster degradation as $\alpha$ rises: generation is fast but produces
    inconsistent abstractions and redundant modules whose maintenance cost
    compounds.
\end{enumerate}

For the comparative statics I suppress the Type-I extreme value shocks
$\varepsilon_t$; with $\varepsilon$ the threshold policy below becomes a
smooth cutoff in choice probabilities and all monotonicity statements apply
to the choice-probability function instead (the logit CCPs inherit the
monotonicity of the choice-specific value difference). Writing the problem
in cost-minimization form, the value function satisfies
\begin{equation}
    V(q) \;=\; \min\Big\{
        \underbrace{c(q,\theta_c) + \beta\, \mathbb{E}\big[V(q') \mid q,\, a=0\big]}_{\text{keep: } K(q)},\;\;
        \underbrace{\kappa + \beta\, \mathbb{E}_\eta\big[V(q_0 + \eta)\big]}_{\text{rationalize: } R}
    \Big\}.
    \label{eq:bellman_costmin}
\end{equation}
The rationalization value $R$ is constant in $q$. Under Assumption A1 below,
$K(q)$ is increasing in $q$, so the optimal policy is a threshold rule:
rationalize if and only if $q \geq q^*$, where
$q^* = \inf\{q : K(q) \geq R\}$ and $\Delta(q) \equiv K(q) - R$ is the
choice-specific value difference.

\paragraph{Assumptions.}
\begin{itemize}
    \item[\textbf{A1}] (\emph{Monotone degradation}) $c(\cdot,\theta_c)$ is
    strictly increasing, and $G_\alpha(\cdot \mid q)$ is stochastically
    increasing in $q$: a more degraded codebase today is (weakly) more
    degraded tomorrow, state by state. This is the discrete-time analogue
    of the convex degradation increment in Equation~\eqref{eq:transition}.
    \item[\textbf{A2}] (\emph{State-dependent acceleration}) $G_\alpha(\cdot
    \mid q)$ is stochastically increasing in $\alpha$, and the shift
    satisfies increasing differences: for every increasing bounded $\phi$
    and $q_H > q_L$,
    \[
        \mathbb{E}_{G_{\alpha}}[\phi \mid q_H] - \mathbb{E}_{G_{\alpha'}}[\phi \mid q_H]
        \;\geq\;
        \mathbb{E}_{G_{\alpha}}[\phi \mid q_L] - \mathbb{E}_{G_{\alpha'}}[\phi \mid q_L],
        \qquad \alpha > \alpha'.
    \]
    AI-induced acceleration is stronger the more degraded the codebase ---
    the compounding (``broken windows'') mechanism of
    \citet{avgeriou2024}, with the convexity evidence of \citet{borg2024}
    and the persistence of AI-introduced issues documented by
    \citet{liu2026}. The post-reset innovation distribution of
    $q_0 + \eta$ is unaffected by $\alpha$ (a fresh codebase carries no
    AI-generated legacy debt at the reset instant).
    \item[\textbf{A3}] (\emph{Regularity}) $c$ bounded on $\mathcal{Q}$,
    $\beta \in (0,1)$, and the threshold is interior for the parameter
    range considered.
\end{itemize}

\begin{proposition}[Cost channel]
\label{prop:cost}
Under A1 and A3, the replacement threshold $q^*(\kappa)$ is strictly
increasing in $\kappa$. Hence AI adoption operating through the cost
channel alone ($\kappa' (\alpha)< 0$, $G$ fixed) strictly lowers the
threshold: firms rationalize at lower degradation levels, and the
per-period rationalization hazard rises.
\end{proposition}

\noindent\emph{Proof sketch.} By the envelope theorem,
$W(q) \equiv \partial V(q)/\partial \kappa$ equals the expected discounted
number of future rationalizations under the optimal policy. The renewal
structure of the problem --- after each rationalization the state
regenerates from $q_0 + \eta$ regardless of history --- yields the identity
$W(q) = m(q)\,\big(1 + B\big)$, where $m(q) \equiv \mathbb{E}_q[\beta^{\tau}]
\in [0,1]$ is the expected discount factor to the \emph{first}
rationalization from state $q$ and
$B \equiv \beta\, \mathbb{E}_\eta[W(q_0+\eta)]$ is common to all states.
Differentiating the value difference,
\[
    \frac{\partial \Delta(q)}{\partial \kappa}
    = \beta\,\mathbb{E}\big[W(q') \mid q, a=0\big] - 1 - B
    = \beta(1+B)\,\mathbb{E}\big[m(q') \mid q, a=0\big] - (1+B)
    \;\leq\; -(1-\beta)(1+B) \;<\; 0,
\]
using $m \leq 1$. The value difference falls strictly and uniformly in
$\kappa$, so the indifference point $q^*$ rises. The full argument,
including the measurable-selection details behind the envelope step, is in
Appendix~\ref{app:proofs}. \hfill$\square$

\begin{proposition}[Transition channel]
\label{prop:transition}
Under A1--A3, the replacement threshold $q^*(\alpha)$ is nonincreasing in
$\alpha$: state-dependent acceleration of degradation lowers the threshold
and raises the rationalization hazard.
\end{proposition}

\noindent\emph{Proof sketch.} The argument is by induction on the value
iteration $V^{n+1} = TV^n$: A1 guarantees each $V^n$ is increasing in $q$;
A2 then implies that the continuation value of \emph{keeping} deteriorates
with $\alpha$ faster at high $q$ than the (reset-anchored) rationalization
value, so the value difference $\Delta(q; \alpha)$ has increasing
differences in $(q, \alpha)$ at every iteration, a property preserved in
the limit \citep[Topkis-style lattice argument;][]{rust1987}. Increasing
differences of $\Delta$ delivers monotonicity of the indifference point.
Details in Appendix~\ref{app:proofs}. \hfill$\square$

\begin{remark}[The sign of the transition channel is state-dependence]
\label{rem:uniform}
A2 is substantive, not technical. If AI accelerated degradation
\emph{uniformly} --- including immediately after a reset --- rationalization
would buy a codebase that itself degrades faster, the option value of
resetting would fall, and the threshold could remain unchanged or even
rise. The transition channel lowers the threshold only insofar as AI-era
degradation \emph{compounds on existing degradation}. This makes A2
directly testable in the adoption design of
Section~\ref{subsec:adoption_design}: the post-adoption shift in
rationalization behavior should be increasing in the repository's
pre-adoption CQI level.
\end{remark}

\subsection{Partial Rationalization: the $(S,s)$ Band and the
Continuous-Maintenance Limit}
\label{subsec:ssband}

The binary model can deliver ``rationalize earlier and more often,'' but it
cannot speak to the \emph{size} of rationalizations: the reset magnitude is
pinned at $q^* - q_0$ by construction. Yet the size margin is where the
data speak loudest (Section~\ref{sec:results}: events arrest degradation
without resetting it). I therefore extend the model to let the firm choose
\emph{how much} to rationalize.

Let the firm pay $\kappa_f + \varphi(z)$ to reduce the state by $z > 0$,
where $\kappa_f = \kappa_f(\alpha)$ is the \emph{fixed} component
(coordination, freeze windows, review and regression risk of any
non-trivial rationalization) and $\varphi$ is an increasing variable cost
with $\varphi(0) = 0$. The binary model is nested: it corresponds to
$\varphi$ forcing a full reset, with $\kappa = \kappa_f + \varphi(q^*-q_0)$.
The central claim of the AI cost channel is then sharper than
$\kappa\!\downarrow$: \emph{AI compresses the fixed component}
$\kappa_f$ --- agents make small and mid-sized rewrites cheap --- which is
exactly the parameter that generates lumpiness in fixed-cost adjustment
models \citep{barro1972, grossmanlaroque1990, caballeroengel1999,
cooperhaltiwanger2006}.

For transparency I derive the deterministic continuous-time case in closed
form; the stochastic case is qualitatively identical
(Remark~\ref{rem:stochastic}). Let the state drift upward at rate
$\mu = \mu(\alpha) > 0$ while operating ($\mu' > 0$: the transition
channel), let flow maintenance cost be $\gamma q$ (the linear term of
Equation~\eqref{eq:cost}; the convex case is in
Appendix~\ref{app:proofs}), let $\varphi(z) = p z$, and impose the natural
floor $q \geq q_0$. In the long-run average-cost problem the optimal policy
is an $(S,s)$ band: let degradation run until $q$ hits the trigger $S$,
then rationalize down to $s = q_0$, i.e.\ jump size $w = S - q_0$. Average
cost per unit time is
\begin{equation}
    \mathrm{AC}(w) \;=\; \underbrace{\frac{\mu\,\big(\kappa_f + p\,w\big)}{w}}_{\text{rationalization}}
    \;+\; \underbrace{\gamma\Big(q_0 + \tfrac{w}{2}\Big)}_{\text{maintenance}},
    \label{eq:ac}
\end{equation}
since cycles have length $w/\mu$ and the state is uniform on $[q_0, q_0+w]$
along a cycle. Minimizing over $w$:
\begin{equation}
    \boxed{\;
    w^* \;=\; \sqrt{\frac{2\,\mu\,\kappa_f}{\gamma}}\,,
    \qquad
    f^* \;=\; \frac{\mu}{w^*} \;=\; \sqrt{\frac{\mu\,\gamma}{2\,\kappa_f}}\,,
    \;}
    \label{eq:band}
\end{equation}
where $w^*$ is the band width (event \emph{amplitude}) and $f^*$ the event
\emph{frequency}.

\begin{proposition}[Band comparative statics and the collapsing-band limit]
\label{prop:band}
In the $(S,s)$ model of Equation~\eqref{eq:ac}:
\begin{enumerate}
    \item[(i)] Frequency rises through \emph{both} channels:
    $\partial f^*/\partial \mu > 0$ and
    $\partial f^*/\partial \kappa_f < 0$.
    \item[(ii)] Amplitude falls through the cost channel but \emph{rises}
    through the transition channel:
    $\partial w^*/\partial \kappa_f > 0$,
    $\partial w^*/\partial \mu > 0$. Under a joint shift
    $(\mu\!\uparrow, \kappa_f\!\downarrow)$, amplitude falls if and only if
    the fixed-cost decline dominates in elasticity:
    $\big|\mathrm{d}\ln \kappa_f\big| > \mathrm{d}\ln \mu$.
    \item[(iii)] The rationalization spend rate,
    $\mu \kappa_f / w^* + p\mu = \sqrt{\mu \gamma \kappa_f / 2} + p\mu$,
    is ambiguous under a joint shift: the transition channel raises it,
    the cost channel lowers it.
    \item[(iv)] As $\kappa_f \to 0$, $w^* \to 0$ and $f^* \to \infty$:
    the policy converges to continuous control at the floor,
    $q \equiv q_0$, with smooth maintenance expenditure at rate $p\,\mu$.
    Rationalization ceases to exist as a discrete event.
\end{enumerate}
\end{proposition}

\noindent The proof is the algebra of Equation~\eqref{eq:band} plus an
elementary limit argument (Appendix~\ref{app:proofs}).

\begin{remark}[Stochastic degradation]
\label{rem:stochastic}
With Brownian noise on the degradation path, the problem is a classical
impulse control of Brownian motion \citep{harrison1983}; the band width
scales as $\kappa_f^{1/3}$ under linear flow cost and, in related quadratic
settings, as $\kappa_f^{1/4}$ \citep{dixit1991}. All qualitative
comparative statics in Proposition~\ref{prop:band} are unchanged.
% TODO(PI): verify the exact exponents and attributions against the primary
% sources before circulation, per project citation policy.
\end{remark}

\paragraph{Predictions and what remains empirical.}
Table~\ref{tab:predictions} collects the signed content of the model.
The joint prediction that is \emph{unambiguous} --- and therefore the
model's primary falsifiable implication --- is that AI adoption raises the
frequency of rationalization activity. Amplitude and the maintenance share
of activity are \emph{signed by which channel dominates}, which is
precisely what the empirical design of
Section~\ref{subsec:adoption_design} is built to measure. Two further
implications discipline the data work. First, the pilot's central pattern
--- friction rising into events and \emph{plateauing} after them, with no
reset to baseline (``rise-then-arrest,''
Figure~\ref{fig:event_study_components}) --- is what an $(S,s)$ firm with
a narrow band looks like in quarterly aggregation: interventions are
frequent and small, so the quarterly state hovers near the trigger rather
than sawtoothing. Read through the model, the pilot pattern is suggestive
evidence that the fixed-cost channel dominates. Second, in the
collapsing-band limit (iv), rationalization events become
\emph{undetectable by construction}: for heavily AI-native codebases the
correct empirical expectation is not larger rewrites but the
disappearance of discrete rewrite events into smooth, elevated
maintenance shares. Event-detection failure at high adoption intensity is
thus a prediction, not a data problem.

\begin{table}[t]
\centering
\caption{Signed predictions of the two-channel $(S,s)$ model}
\label{tab:predictions}
\begin{tabular}{lccc}
\hline
 & Cost channel & Transition channel & Joint shift \\
Object & ($\kappa_f \downarrow$) & ($\mu \uparrow$) & ($\kappa_f\downarrow,\ \mu\uparrow$) \\
\hline
Rationalization frequency $f^*$ & $\uparrow$ & $\uparrow$ & $\uparrow$ (unambiguous) \\
Event amplitude $w^*$ & $\downarrow$ & $\uparrow$ & $\downarrow$ iff $|\mathrm{d}\ln\kappa_f| > \mathrm{d}\ln\mu$ \\
Replacement trigger $S^*$ & $\downarrow$ & $\uparrow$ & ambiguous \\
Rationalization spend rate & $\downarrow$ & $\uparrow$ & ambiguous (empirical) \\
Threshold effect by baseline $q$ & --- & stronger at high $q$ (A2) & testable heterogeneity \\
\hline
\end{tabular}
\end{table}

\paragraph{Implications for estimation.}
The $(S,s)$ structure changes what the data identify. In the binary model,
the CCP estimator recovers $(\theta_1, \theta_2, \kappa)$ from
rationalization hazards alone. In the band model, the \emph{joint}
distribution of event frequency and event amplitude separately identifies
the fixed cost $\kappa_f$ (from amplitude, via Equation~\eqref{eq:band})
and the drift $\mu$ (from frequency given amplitude) --- so the two AI
channels are separately identified from moments the panel already
measures: refactor-episode sizes (churn and commit shares during event
windows) and inter-event durations. This is the estimation target for the
scaled sample; the pilot's role is to validate the event and amplitude
measurement underlying those moments.

% =========================================================================
% APPENDIX MATERIAL — move to the paper's appendix
% =========================================================================

\section*{Appendix: Proofs for Sections~\ref{subsec:twochannel}--\ref{subsec:ssband}}
\label{app:proofs}

\paragraph{Proof of Proposition~\ref{prop:cost}.}
Fix the primitives except $\kappa$. Standard arguments (bounded costs,
$\beta < 1$) give existence and uniqueness of $V$ and, under A1, that $V$
and $K$ are increasing in $q$, so the optimal policy is the threshold rule
$a = \mathbf{1}\{q \geq q^*\}$ with $q^* = \inf\{q : \Delta(q) \geq 0\}$,
$\Delta = K - R$.

\emph{Step 1 (envelope).} Let $\Pi$ denote the optimal policy at $\kappa$.
For fixed policy $\Pi$, the value is affine in $\kappa$:
$V_\Pi(q; \kappa) = A_\Pi(q) + \kappa\, W_\Pi(q)$, where
$W_\Pi(q) = \mathbb{E}^\Pi_q\big[\sum_{i \geq 1} \beta^{\tau_i}\big]$ is
the expected discounted number of rationalizations at times
$\tau_1 < \tau_2 < \cdots$. By the envelope theorem for dynamic
discrete-choice values (or directly: two-sided policy-suboptimality
bounds), $\partial V/\partial \kappa = W_{\Pi}(q)$ at almost every
$\kappa$.

\emph{Step 2 (renewal identity).} At $\tau_1$ the firm pays and the state
regenerates from $q_0 + \eta$, independently of history. Hence
\[
    W(q) \;=\; \mathbb{E}_q\big[\beta^{\tau_1}\big]\Big(1 +
    \beta\,\mathbb{E}_\eta\big[W(q_0+\eta)\big]\Big)
    \;=\; m(q)\,(1 + B),
\]
with $m(q) = \mathbb{E}_q[\beta^{\tau_1}] \in [0, 1]$ and
$B = \beta\,\mathbb{E}_\eta[W(q_0+\eta)] \geq 0$ independent of $q$.

\emph{Step 3 (uniform strict decrease of $\Delta$).}
\[
    \frac{\partial \Delta(q)}{\partial \kappa}
    = \beta\,\mathbb{E}\big[W(q') \mid q, 0\big] - \Big(1 +
      \beta\,\mathbb{E}_\eta\big[W(q_0+\eta)\big]\Big)
    = \beta (1+B)\, \mathbb{E}\big[m(q') \mid q, 0\big] - (1 + B).
\]
Since $m \leq 1$, the right side is at most
$\beta(1+B) - (1+B) = -(1-\beta)(1+B) < 0$. Thus for $\kappa_1 > \kappa_0$,
$\Delta(q; \kappa_1) < \Delta(q; \kappa_0)$ for every $q$, and therefore
$q^*(\kappa_1) \geq q^*(\kappa_0)$, strictly when interior (A3).
\hfill$\blacksquare$

\paragraph{Proof of Proposition~\ref{prop:transition} (induction).}
Let $T_\alpha$ denote the Bellman operator at intensity $\alpha$ and
$V^0 \equiv 0$. We show by induction that each iterate satisfies
(P$_n$): $V^n(\cdot\,; \alpha)$ is increasing in $q$, and
$D^n(q) \equiv V^n(q; \alpha) - V^n(q; \alpha')$ is nonnegative and
increasing in $q$ for $\alpha > \alpha'$.

\emph{Base.} Trivial.

\emph{Step.} Monotonicity in $q$: $c$ increasing (A1) and $G_\alpha$
stochastically increasing in $q$ (A1) imply $K^{n+1}$ increasing; $R^{n+1}$
constant; the min of an increasing and a constant function is increasing.

Nonnegativity of $D^{n+1}$: by A2 (FOSD in $\alpha$) and $V^n$ increasing,
$\mathbb{E}_{G_\alpha}[V^n \mid q] \geq \mathbb{E}_{G_{\alpha'}}[V^n \mid q]$;
adding $D^n \geq 0$ inside the expectations preserves the ranking; the
reset branch is unchanged in $\alpha$ except through
$\mathbb{E}_\eta[D^n(q_0+\eta)] \geq 0$. The pointwise min of two functions
each larger at $\alpha$ is larger at $\alpha$.

Monotonicity of $D^{n+1}$ in $q$: on the keep region,
\[
    D^{n+1}(q) = \beta\Big(
      \underbrace{\mathbb{E}_{G_\alpha}[V^n(\cdot;\alpha) \mid q]
      - \mathbb{E}_{G_{\alpha'}}[V^n(\cdot;\alpha) \mid q]}_{\text{direct shift}}
      + \underbrace{\mathbb{E}_{G_{\alpha'}}[D^n \mid q]}_{\text{inherited}}
    \Big).
\]
The direct-shift term is increasing in $q$ by A2 (increasing differences)
applied to the increasing function $V^n(\cdot;\alpha)$; the inherited term
is increasing in $q$ because $D^n$ is increasing (induction) and
$G_{\alpha'}$ is stochastically increasing in $q$ (A1). On the rationalize
region $D^{n+1}$ is constant at
$\beta\,\mathbb{E}_\eta[D^n(q_0+\eta)]$, which is (weakly) below its value
on the keep region at the boundary; the piecewise structure preserves
monotonicity across the boundary because the keep branch attains the min
exactly where its (larger) $\alpha$-sensitivity applies.

(P$_n$) passes to the limit $V = \lim V^n$ (uniform convergence,
$\beta$-contraction). Increasingness of $D$ in $q$ is precisely increasing
differences of $-\Delta$...

Concretely: $\Delta(q; \alpha) - \Delta(q; \alpha') =
\beta\big(\mathbb{E}_{G_\alpha}[V(\cdot;\alpha)\mid q] -
\mathbb{E}_{G_{\alpha'}}[V(\cdot;\alpha')\mid q]\big) -
\beta\,\mathbb{E}_\eta[D(q_0+\eta)]$. The first bracket equals the direct
shift plus inherited term above evaluated at the limit, which is
increasing in $q$ and, at the reset-anchored comparison point, at least
$\mathbb{E}_\eta[D(q_0+\eta)]$ for all $q$ in a neighborhood of the
threshold provided $q_0 + \eta \leq q^*$ a.s.\ (resets land strictly
inside the continuation region --- an implication of A3 interiority).
Hence $\Delta(\cdot; \alpha) \geq \Delta(\cdot; \alpha')$ near the
threshold, and $q^*(\alpha) \leq q^*(\alpha')$.
\hfill$\blacksquare$

% NOTE(PI): the last paragraph is the step a careful referee will probe —
% the comparison of the keep-branch alpha-sensitivity against the
% reset-branch alpha-sensitivity near the threshold. The sufficient
% condition doing the work is A2 (state-dependent acceleration) together
% with resets landing inside the continuation region. If you want a
% cleaner presentation, an alternative is to assume directly that
% E[V|q,0] - E[V|q_0-reset] has increasing differences in (q, alpha) and
% relegate primitives to a lemma. Please scrutinize line by line.

\paragraph{Proof of Proposition~\ref{prop:band}.}
(i)--(iii) are direct differentiation of Equation~\eqref{eq:band}:
$\ln w^* = \tfrac12(\ln 2 + \ln\mu + \ln\kappa_f - \ln\gamma)$, so
$\mathrm{d}\ln w^* = \tfrac12(\mathrm{d}\ln\mu + \mathrm{d}\ln\kappa_f)$,
negative iff $|\mathrm{d}\ln\kappa_f| > \mathrm{d}\ln\mu$; and
$\ln f^* = \tfrac12(\ln\mu + \ln\gamma - \ln 2 - \ln\kappa_f)$ is
increasing in $\mu$ and decreasing in $\kappa_f$. The spend rate is
$\sqrt{\mu\gamma\kappa_f/2} + p\mu$, increasing in $\mu$ and increasing in
$\kappa_f$ (hence falling as AI compresses $\kappa_f$). For (iv), as
$\kappa_f \to 0$ the minimized average cost converges to
$\gamma q_0 + p\mu$, the cost of holding $q$ at the floor while paying the
variable cost of absorbing degradation flow $\mu$ continuously; the band
$[q_0, q_0 + w^*]$ collapses to the point $q_0$ and cycle length
$w^*/\mu \to 0$.

\emph{Convex flow cost.} With $c(q) = \theta_1 q + \theta_2 q^2$ the
cycle-average maintenance cost is
$\theta_1(q_0 + w/2) + \theta_2\big(q_0^2 + q_0 w + w^2/3\big)$ and the
first-order condition becomes
$\mu\kappa_f / w^2 = \theta_1/2 + \theta_2 q_0 + \tfrac{2}{3}\theta_2 w$,
which has a unique positive root with the same limits and the same signs
of all comparative statics; convexity ($\theta_2 > 0$) shrinks $w^*$
relative to the linear case, reinforcing (iv).
\hfill$\blacksquare$
